% ===========================================================================
% Rapport ProSys -- Odoo 14 en haute disponibilite sur LXC
%
% Compilation (depuis la racine du depot) :
%     tectonic rapport-odoo14.tex
%
% Les schemas sont lus dans docs/schemas-patroni-odoo14/ (\graphicspath).
% Leur numero dans le rapport (Schema 1 a 14) est celui de leur fichier
% (01-... a 14-...) : les sections suivent l'ordre du catalogue.
% ===========================================================================
\documentclass[11pt,a4paper]{article}

\usepackage{fontspec}
\usepackage[french]{babel}
\usepackage[a4paper,margin=1.9cm,headheight=14pt,includehead,includefoot]{geometry}
\usepackage{graphicx}
\usepackage{float}
\usepackage{array}
\usepackage{tabularx}
\usepackage{booktabs}
\usepackage[table]{xcolor}
\usepackage{enumitem}
\usepackage{listings}
\usepackage{tcolorbox}
\tcbuselibrary{breakable,skins}
\usepackage{caption}
\usepackage{fancyhdr}
\usepackage[hidelinks]{hyperref}

% --- Chemins et reglages generaux ------------------------------------------
% Schemas 1 a 14 par leur seul nom ; schemas 15 a 24 par leur sous-dossier
% de docs/ (schemas-lxc/, schemas-glusterfs/, schemas-architecture-ha/).
\graphicspath{{docs/schemas-patroni-odoo14/}{docs/}}

% Nom de fichier dans un titre de section (signet PDF sans mise en forme).
\newcommand{\pb}[1]{\texorpdfstring{\texttt{#1}}{#1}}
\setlength{\emergencystretch}{2em}
\setlist{itemsep=2pt,topsep=4pt}

% Code en ligne : underscores et chemins sans echappement.
% Ne pas l'utiliser dans un titre ni une legende.
\newcommand{\code}[1]{\texttt{\detokenize{#1}}}

\newcolumntype{L}{>{\raggedright\arraybackslash}X}
\newcolumntype{P}[1]{>{\raggedright\arraybackslash}p{#1}}

\addto\captionsfrench{%
  \renewcommand{\figurename}{Schéma}%
  \renewcommand{\listfigurename}{Liste des schémas}%
  \renewcommand{\tablename}{Tableau}%
}
\captionsetup{font=small,labelfont=bf}

% --- Blocs de commandes ------------------------------------------------------
\lstset{
  basicstyle=\ttfamily\small,
  breaklines=true,
  backgroundcolor=\color{gray!8},
  frame=single,
  rulecolor=\color{gray!40},
  commentstyle=\color{gray!70!black},
  showstringspaces=false,
  keepspaces=true,
  columns=flexible,
  language=bash,
}

% --- Schema et lecture du schema --------------------------------------------
% \schema{fichier}{titre}
\newcommand{\schema}[2]{%
  \begin{figure}[H]
    \centering
    \includegraphics[width=\textwidth]{#1}
    \caption{#2}
    \label{fig:#1}
  \end{figure}}

% Encadre place sous chaque schema : resume, points de lecture, sources.
\newenvironment{lecture}[1]{%
  \begin{tcolorbox}[breakable, enhanced, colback=gray!4, colframe=gray!55,
    boxrule=0.5pt, arc=1mm, left=2mm, right=2mm, top=1mm, bottom=1mm,
    fonttitle=\bfseries\small, title={Lecture du schéma~\thefigure}]
  % Fer a gauche : les longs noms de fichiers ne peuvent pas etre coupes.
  \raggedright\emph{#1}\par\smallskip\small}
  {\end{tcolorbox}}

% --- Entetes et pieds de page -----------------------------------------------
\pagestyle{fancy}
\fancyhf{}
\fancyhead[L]{\small ProSys -- Odoo 14 en haute disponibilité}
\fancyhead[R]{\small\nouppercase{\leftmark}}
\fancyfoot[L]{\small ADMIN}
\fancyfoot[R]{\small \thepage}
\renewcommand{\headrulewidth}{0.4pt}

\hypersetup{
  % La page de titre remet le compteur de pages a 1 : sans cela, deux pages
  % portent le meme nom de destination PDF ("@page.1 already defined").
  hypertexnames=false,
  pdftitle={ProSys - Odoo 14 en haute disponibilité sur LXC},
  pdfauthor={ADMIN},
}

% ===========================================================================
\begin{document}

% --- Page de titre -----------------------------------------------------------
\begin{titlepage}
  \centering
  \vspace*{1.5cm}
  {\large\scshape Rapport de projet\par}
  \vspace{0.8cm}
  {\Huge\bfseries ProSys -- Odoo 14\\[4pt] en haute disponibilité sur LXC\par}
  \vspace{0.6cm}
  {\Large Architecture, déploiement Ansible et exploitation\par}
  \vfill
  {\large ADMIN\par}
  \vspace{0.2cm}
 
\end{titlepage}

% --- Resume -----------------------------------------------------------------
\begin{abstract}
Ce rapport décrit le projet ProSys : le déploiement par Ansible de trois
instances \textbf{Odoo~14} qui partagent une base \textbf{PostgreSQL~15}
rendue hautement disponible par \textbf{Patroni}, \textbf{etcd} et
\textbf{HAProxy}, et un filestore commun porté par un volume
\textbf{GlusterFS}. Chaque instance tourne dans un conteneur LXC Debian~12,
posé sur son propre hôte et relié au réseau en IPVLAN.

Le rapport présente l'architecture, les mécanismes de haute disponibilité
(réplication, bascule automatique et planifiée, quorum, réintégration), la
sécurité des flux, la sauvegarde, l'installation et l'exploitation, puis
détaille chacun des huit playbooks. Les vingt-quatre schémas sont suivis
chacun d'une lecture commentée : quatorze décrivent l'architecture actuelle,
dix proviennent des éditions précédentes et montrent son évolution. Les
limites connues et les points restant à compléter sont signalés.


\end{abstract}

\tableofcontents
\clearpage

% ===========================================================================
\section{Introduction}
% ===========================================================================

\subsection{Contexte et objectifs}

ProSys exploite Odoo~14 pour plusieurs utilisateurs. Le projet répond à trois
objectifs :

\begin{itemize}
  \item \textbf{servir une base unique} à trois instances Odoo, pour que toutes
        voient les mêmes données ;
  \item \textbf{rendre la base hautement disponible} : si le serveur
        PostgreSQL primaire tombe, une réserve prend sa place sans intervention ;
  \item \textbf{rendre le déploiement reproductible} : toute l'installation est
        décrite par des playbooks Ansible versionnés.
\end{itemize}

\subsection{Principes SRE appliqués}

L'ingénierie de la fiabilité (SRE, \emph{Site Reliability Engineering}) traite
l'exploitation comme un problème d'ingénierie : la fiabilité se définit et se
mesure, le travail répétitif s'automatise, les changements sont rendus sûrs, et
chaque incident sert à renforcer le système. Le projet applique ces principes
de la façon suivante ; les points encore ouverts sont signalés.

\begin{itemize}
  \item \textbf{Supprimer les points de défaillance uniques.} Trois serveurs
        PostgreSQL sous Patroni, un quorum etcd de 2 sur 3, un HAProxy par
        serveur et trois instances Odoo : la perte d'un nœud de données est
        tolérée sans intervention. Restent uniques : le cron d'Odoo, porté par
        odoo1 ; le serveur 192.168.36.8, dont dépend le montage du volume au
        démarrage, aucun serveur de secours n'étant déclaré ; l'hôte de
        sauvegarde.
  \item \textbf{Automatiser plutôt que répéter} (réduction du \emph{toil}).
        Tout le déploiement est décrit par des playbooks Ansible versionnés.
        Ils sont idempotents : rejoués, ils remettent en conformité sans rien
        casser, et \code{--check --diff} montre ce qui changerait avant d'agir.
  \item \textbf{Rendre les changements sûrs.} Les opérations risquées sont
        progressives, pour limiter le rayon d'impact : montage du filestore un
        hôte à la fois, relance de Patroni un serveur à la fois et cluster en
        pause, redémarrages des réserves avant le primaire. À l'inverse, une
        reprise à moitié faite étant pire qu'une reprise pas commencée, un
        échec sur un serveur PostgreSQL arrête les trois
        (\code{any_errors_fatal}).
  \item \textbf{Échouer de façon sûre} (garde-fous). Odoo ne démarre pas sans
        son volume partagé, et le rôle \code{odoo14} refuse de configurer une
        instance sans lui. Le rôle \code{pg_ha} refuse de créer une nouvelle
        autorité de certification et s'arrête sur « refus » avant toute
        modification. La base n'est initialisée que par odoo1.
  \item \textbf{Apprendre des incidents.} Les incidents rencontrés sont
        documentés dans le code, avec leur date et leur cause, et chacun a
        donné un garde-fou : sauvegarde refusée par la base faute de règle pour
        son hôte (24/09/2026), Odoo écrivant sur son disque local faute de
        volume (25/09/2026), instances arrêtées par un mauvais chemin du
        filestore dans le volume (29 et 30/09/2026), base voyant les Odoo sous
        l'adresse du routeur NAT (20/10/2026). C'est l'esprit du post-mortem
        sans recherche de coupable : corriger le système, pas les personnes.
  \item \textbf{Rendre l'environnement reproductible.} Les versions des
        collections Ansible sont épinglées, et un contrôle optionnel (tag
        \code{preflight}) les vérifie avant tout contact avec les hôtes.
        L'archive de Python~3.8.20 est contrôlée par son empreinte SHA256 ; les
        dépendances Python viennent d'un fichier de versions figées lorsqu'il
        existe.
  \item \textbf{Prévoir la capacité.} La règle du projet borne les connexions
        à la base : 3 nœuds $\times$ (2 workers + 1 cron) $\times$ 16
        connexions, soit 144 au plus, pour 200 autorisées par
        \code{max_connections}.
  \item \textbf{Observer et vérifier.} Chaque déploiement se termine par des
        contrôles actifs : un primaire et deux réserves, retard inférieur à
        16~Mo, vraie connexion TLS à travers chaque HAProxy, montage GlusterFS
        dans chaque conteneur. L'état se lit par \code{patronictl list}, l'API
        Patroni, la page d'état HAProxy et les journaux systemd.
        \emph{À compléter :} le dépôt ne déploie ni supervision continue ni
        alertes.
  \item \textbf{Mesurer la fiabilité} (SLI, SLO, budget d'erreur).
        \emph{À définir :} le dépôt ne fixe pas encore d'objectifs chiffrés.
        Indicateurs proposés : disponibilité de \code{/web/login} sur chaque
        instance, disponibilité des écritures SQL, temps de rétablissement
        après une bascule (RTO) et perte de données admise (RPO). Les réglages
        comme \code{ttl=30 s} ou \code{maximum_lag_on_failover} ne sont pas
        des mesures : seuls des tests de bascule et de restauration
        chronométrés les fourniront.
\end{itemize}

Le rapport décrit la configuration prévue par le code ; l'état réel des
serveurs, les délais de bascule et la restauration restent à constater en
exploitation.

% ===========================================================================
\section{Architecture générale}
% ===========================================================================

\subsection{Vue d'ensemble}

L'architecture sépare une zone applicative, où tournent les trois Odoo, et une
zone de données, où tournent les trois serveurs PostgreSQL. Le filestore est
porté par un volume GlusterFS extérieur au projet.

\schema{01-architecture-globale-ha}{Architecture globale : Odoo et PostgreSQL HA}

\begin{lecture}{Trois instances applicatives, trois nœuds de données et un filestore partagé.}
\begin{itemize}
  \item \textbf{Zone application} (\code{192.168.36.0/24}) : les utilisateurs
        accèdent directement à l'une des trois instances Odoo, en HTTP sur
        8069 et sur 8072 pour le bus de notifications. Les trois instances
        tournent dans \code{/opt/odoo14-venv} avec Python~3.8 ; odoo1 porte en
        plus l'unique cron.
  \item \textbf{Routage L3} : un routeur externe relie les deux réseaux et
        applique un SNAT vers \code{192.168.33.254}.
  \item \textbf{Zone données} (\code{192.168.33.0/24}) : chaque nœud porte
        PostgreSQL~15 (5432), Patroni (8008), etcd (2379 et 2380) et HAProxy
        (5000).
  \item \textbf{Accès SQL} : Odoo connaît les trois points d'entrée HAProxy
        sur TCP/5000 ; libpq choisit l'adresse disponible, et le trafic SQL
        est chiffré en TLS.
  \item \textbf{Rôles} : seul le primaire reçoit le SQL. Les répliques ne
        reçoivent que les WAL, le journal des transactions du primaire
        (expliqué à la section~\ref{sec:replication}), de façon asynchrone. Patroni et etcd
        coordonnent le verrou du leader avec un quorum de 2 sur 3.
  \item \textbf{Filestore} : le volume GlusterFS externe \code{sre} est monté
        directement dans chaque conteneur Odoo, sur \code{/var/lib/odoo}.
        \code{192.168.36.8:24007} fournit la description du volume, puis le
        client dialogue avec les serveurs de briques.
  \item Le placement « odoodb1 primaire » n'est qu'un exemple : le primaire
        peut être n'importe lequel des trois nœuds.
\end{itemize}
\end{lecture}

\subsection{Machines et conteneurs}

Chaque hôte Debian~12 porte un seul conteneur LXC. Le placement est défini
dans \code{inventory.yaml} et dans \code{host_vars/}.

\begin{table}[H]
\centering\small
\begin{tabularx}{\textwidth}{L P{2cm} P{2.6cm} P{2.4cm} P{1.7cm}}
\toprule
\rowcolor{gray!15}
\textbf{Rôle} & \textbf{Conteneur} & \textbf{IP du conteneur} & \textbf{Hôte} & \textbf{Carte} \\
\midrule
odoo1 -- Odoo, initialisation de la base, cron & \code{devodoo14} & 192.168.36.161 & 192.168.36.21 & \code{enp5s0} \\
odoo2 -- Odoo                                   & \code{devodoo14} & 192.168.36.162 & 192.168.36.22 & \code{enp5s0} \\
odoo3 -- Odoo                                   & \code{devodoo14} & 192.168.36.163 & 192.168.36.23 & \code{enp5s0} \\
odoodb1 -- PostgreSQL, Patroni, etcd, HAProxy   & \code{devpg15}   & 192.168.33.161 & 192.168.33.21 & \code{enp1s0} \\
odoodb2 -- PostgreSQL, Patroni, etcd, HAProxy   & \code{devpg15}   & 192.168.33.162 & 192.168.33.22 & \code{enp1s0} \\
odoodb3 -- PostgreSQL, Patroni, etcd, HAProxy   & \code{devpg15}   & 192.168.33.163 & 192.168.33.23 & \code{enp1s0} \\
\bottomrule
\end{tabularx}
\caption{Hôtes et conteneurs déclarés dans l'inventaire}
\end{table}

\begin{table}[H]
\centering\small
\begin{tabularx}{\textwidth}{P{4.6cm} P{4.2cm} L}
\toprule
\rowcolor{gray!15}
\textbf{Élément extérieur} & \textbf{Adresse} & \textbf{Usage} \\
\midrule
Volume GlusterFS répliqué \code{sre} & 192.168.36.8 (description du volume) & Filestore et sessions communs aux trois Odoo \\
Routeur entre les deux réseaux & passerelles \code{.254} & NAT : la base voit les Odoo sous 192.168.33.254 \\
Hôte de sauvegarde & 192.168.36.21 (\code{backup_host}) & \code{pg_dump} et archive du filestore dans \code{/srv/sauvegardes} \\
\bottomrule
\end{tabularx}
\caption{Éléments extérieurs que le déploiement suppose en place}
\end{table}

\schema{02-deploiement-hotes-lxc}{Déploiement : hôtes Debian et conteneurs LXC}

\begin{lecture}{Six hôtes déclarés ; chaque conteneur possède une adresse LAN via IPVLAN L2.}
\begin{itemize}
  \item Chaque boîte représente un hôte Debian~12 et le conteneur qu'il porte.
        Les hôtes 192.168.36.21 à 23 utilisent la carte \code{enp5s0} et
        portent odoo1 à odoo3. Les hôtes 192.168.33.21 à 23 utilisent
        \code{enp1s0} et portent odoodb1 à odoodb3.
  \item Les conteneurs sont reliés en \textbf{IPVLAN L2} : ils partagent
        l'adresse MAC de leur hôte, la seule que le réseau accepte.
  \item \textbf{Conteneurs Odoo} : Debian~12, Odoo~14, Python~3.8, venv
        local \code{/opt/odoo14-venv}, volume GlusterFS monté dans le
        conteneur.
  \item \textbf{Conteneurs PostgreSQL} : PostgreSQL~15 avec Patroni, etcd et
        HAProxy. Les données PostgreSQL sont locales à chaque nœud et
        recopiées par réplication.
  \item Le routage externe relie 192.168.36.0/24 et 192.168.33.0/24 ; Odoo
        atteint les HAProxy sur le port 5000.
  \item \textbf{Domaines de panne} : les hôtes sont distincts dans
        l'inventaire, mais leur placement physique n'est pas décrit. En
        IPVLAN, un hôte ne joint pas son propre conteneur sans chemin adapté.
\end{itemize}
\end{lecture}

\subsection{Filestore commun}

Chaque conteneur Odoo monte \emph{lui-même} \code{192.168.36.8:/sre} sur
\code{/var/lib/odoo}. Le volume contient \code{filestore/prod} et
\code{sessions}. Trois garde-fous protègent ce montage :

\begin{itemize}
  \item l'unité systemd d'Odoo exige le montage (\code{RequiresMountsFor}).
        Sans volume, Odoo ne démarre pas, au lieu d'écrire sur le disque
        local ;
  \item le rôle \code{odoo14} refuse de configurer une instance dont
        \code{/var/lib/odoo} n'est pas un montage GlusterFS ;
  \item le compte \code{odoo} a le même UID et le même GID (105/112) sur les
        trois conteneurs. Sur un volume partagé, ce sont les numéros qui
        comptent.
\end{itemize}

L'ancien montage sur l'hôte suivi d'un bind vers le conteneur n'existe plus.
Seul l'hôte de sauvegarde garde \code{/mnt/odoo}, en lecture seule.

% ===========================================================================
\section{Haute disponibilité de PostgreSQL}
% ===========================================================================

\subsection{Rôle des composants}

La haute disponibilité repose sur trois composants installés dans chacun des
trois conteneurs PostgreSQL. Patroni pilote le PostgreSQL local, etcd conserve
l'état du cluster, et HAProxy dirige les connexions vers le primaire.

\schema{03-orchestration-patroni}{Patroni orchestre le rôle de chaque PostgreSQL}

\begin{lecture}{Les données circulent entre PostgreSQL ; les décisions de rôle passent par Patroni et etcd.}
\begin{itemize}
  \item Odoo, par libpq, connaît trois adresses HAProxy. Le SQL chiffré passe
        par l'un des trois HAProxy (port 5000), puis rejoint le primaire sur
        5432.
  \item HAProxy contrôle le rôle de chaque serveur par \code{GET /primary}
        auprès des trois agents Patroni : la réponse 200 désigne le primaire
        qui détient le verrou.
  \item Les WAL vont directement du primaire vers les répliques ; ils ne
        traversent pas Patroni.
  \item \textbf{Patroni} (API 8008) : un agent indépendant par nœud, qui
        démarre, arrête, configure et promeut son PostgreSQL local.
  \item \textbf{etcd} (quorum 2/3) : conserve le verrou du leader et la
        configuration ; clients en HTTPS sur 2379, échanges entre membres sur
        2380.
  \item En résumé : Patroni supervise et coordonne, HAProxy route le TCP,
        PostgreSQL exécute et réplique, etcd conserve l'état de coordination.
\end{itemize}
\end{lecture}

\subsection{Connexion d'Odoo à la base}

Odoo reçoit dans \code{/etc/odoo/odoo.conf} la liste des trois serveurs et le
port des HAProxy :

\begin{lstlisting}
db_host = 192.168.33.161,192.168.33.162,192.168.33.163
db_port = 5000
\end{lstlisting}

Comme le routeur fait du NAT, la base reconnaît chaque instance à son
\textbf{certificat client} et non à son adresse.

\schema{04-connexion-odoo-haproxy}{Connexion Odoo : du venv au primaire courant}

\begin{lecture}{libpq essaie les points d'entrée dans l'ordre et garde une connexion vers un serveur qui accepte les écritures.}
\begin{itemize}
  \item Odoo~14 s'exécute dans le venv et passe par psycopg2 puis libpq, avec
        \code{PGTARGETSESSIONATTRS=read-write} et
        \code{PGSSLMODE=verify-full}.
  \item libpq essaie les trois adresses, port 5000, et ouvre une nouvelle
        connexion vers la première disponible. Le choix du proxy dépend de sa
        disponibilité ; c'est HAProxy qui sélectionne le primaire.
  \item Chaque HAProxy (sur odoodb1, odoodb2 et odoodb3) est un relais TCP
        transparent qui sonde les trois agents Patroni sur 8008.
  \item Toutes les connexions aboutissent au \textbf{primaire courant}
        (5432), qui reçoit toutes les lectures et écritures d'Odoo. Le
        certificat du serveur est vérifié, et celui d'Odoo est exigé.
  \item Le TLS PostgreSQL va de bout en bout : HAProxy ne le termine pas.
  \item Après une bascule, HAProxy ferme les sessions attachées à la cible
        retirée. Les clients doivent se reconnecter, et les transactions en
        cours ne sont pas transférées.
\end{itemize}
\end{lecture}

\subsection{Réplication}
\label{sec:replication}

\paragraph{Qu'est-ce que le WAL ?} Le WAL (\emph{Write-Ahead Log}, « journal
écrit en premier ») est le journal des transactions de PostgreSQL. Avant de
modifier ses fichiers de données, PostgreSQL écrit chaque changement dans ce
journal, sous forme d'enregistrements numérotés à la suite, et le force sur le
disque au moment du \code{COMMIT}. Le journal est découpé en fichiers de
16~Mo, rangés dans le répertoire \code{pg_wal/}. On parle couramment « des
WAL » pour désigner ces fichiers ou les enregistrements qu'ils contiennent.

Ce principe sert à trois choses dans le projet :
\begin{itemize}
  \item \textbf{reprise après une panne} : au redémarrage, PostgreSQL rejoue
        le WAL pour retrouver un état cohérent, même si les fichiers de
        données n'avaient pas encore été mis à jour ;
  \item \textbf{réplication} : le primaire envoie son WAL en continu
        (\emph{streaming}) aux répliques, qui le rejouent. Elles obtiennent
        ainsi une copie exacte de la base. Ce ne sont pas les requêtes SQL qui
        sont transmises, mais les modifications physiques des pages de
        données ;
  \item \textbf{réintégration d'un ancien primaire} : \code{pg_rewind}
        s'appuie sur l'historique du WAL pour ne recopier que les blocs qui
        ont divergé. C'est pour cela que \code{wal_log_hints = on} est activé.
\end{itemize}

Quatre notions reviennent dans la suite :
\begin{itemize}
  \item \textbf{\code{wal_level = replica}} : le WAL contient assez
        d'informations pour alimenter des répliques ;
  \item \textbf{slot de réplication} : le primaire conserve le WAL qu'une
        réplique n'a pas encore reçu. \code{max_slot_wal_keep_size = 2GB} borne
        cette réserve, pour qu'une réplique absente ne remplisse pas le disque
        du primaire ;
  \item \textbf{retard} (\emph{lag}) : la quantité de WAL produite par le
        primaire et pas encore reçue ou rejouée par une réplique, mesurée en
        octets. C'est elle que comparent \code{maximum_lag_on_failover}
        (1~Mio) et le contrôle final d'Ansible (16~Mo) ;
  \item \textbf{asynchrone} : le primaire confirme le \code{COMMIT} à Odoo sans
        attendre qu'une réplique ait reçu le WAL correspondant. C'est rapide,
        mais les toutes dernières transactions peuvent manquer sur une réplique
        promue après une panne.
\end{itemize}

\schema{05-replication-postgresql}{Réplication : un primaire et deux répliques}

\begin{lecture}{La réplication physique des WAL reste asynchrone ; le quorum etcd ne vaut pas acquittement des écritures.}
\begin{itemize}
  \item Le \textbf{primaire} (exemple : odoodb1) traite les lectures et
        écritures d'Odoo, produit les WAL et utilise des slots physiques
        gérés par Patroni.
  \item Les \textbf{répliques} odoodb2 et odoodb3 reçoivent puis rejouent les
        WAL, en lecture seule, par une connexion de réplication TLS sur 5432.
        Elles ne reçoivent aucune requête d'Odoo.
  \item \textbf{Rétention} : \code{use_slots = true},
        \code{wal_level = replica}, \code{max_wal_senders = 10},
        \code{max_replication_slots = 10},
        \code{max_slot_wal_keep_size = 2GB}.
  \item \textbf{Éligibilité} : \code{maximum_lag_on_failover} vaut
        1\,048\,576 octets (1~Mio). Le contrôle final d'Ansible exige un
        retard inférieur ou égal à 16~Mo. Les deux seuils répondent à des
        contrôles différents.
  \item \textbf{Perte possible} : une écriture confirmée sur le primaire peut
        manquer sur une réplique promue. Le seuil de 1~Mio ne garantit pas une
        perte maximale de 1~Mio.
\end{itemize}
\end{lecture}

\subsection{Bascule automatique}

\schema{06-bascule-automatique}{Bascule automatique après panne du primaire}

\begin{lecture}{Deux mécanismes avancent en parallèle : la coordination Patroni/etcd et le retrait de la cible par HAProxy.}
\textbf{Coordination, en quatre étapes :}
\begin{enumerate}
  \item \emph{Primaire perdu} : panne de PostgreSQL, du conteneur, ou de la
        communication du leader.
  \item \emph{Verrou indisponible} : le verrou expire ou se perd, alors que
        le quorum etcd reste disponible.
  \item \emph{Réplique éligible} : une réplique dont l'état et le retard sont
        compatibles acquiert le verrou.
  \item \emph{Promotion} : elle devient le nouveau primaire (exemple :
        odoodb2).
\end{enumerate}
\textbf{Connexions :} les sondes \code{/primary} échouent et HAProxy retire
la cible, en fermant les sessions qui lui étaient liées. Dès que le nouveau
primaire répond 200, HAProxy lui route les connexions SQL. Odoo ouvre alors
une nouvelle session ; les transactions en cours sont à traiter par
l'application.

\textbf{Réglages déclarés :} \code{ttl=30 s}, \code{loop_wait=10 s},
\code{retry_timeout=10 s} ; sondes HAProxy \code{inter=3 s}, \code{fall=3},
\code{rise=2}. Ces valeurs ne constituent pas un temps de rétablissement
mesuré. Une bascule ne transfère pas les transactions en cours et, en
réplication asynchrone, ne garantit pas l'absence de perte.
\end{lecture}

\subsection{Quorum et partitions réseau}

\schema{07-quorum-et-partitions}{Quorum etcd et partitions réseau}

\begin{lecture}{La majorité autorise la coordination ; Patroni utilise le verrou pour gérer le rôle de PostgreSQL.}
\begin{itemize}
  \item \textbf{A -- Trois membres joignables} : la majorité (au moins 2 sur
        3) est atteinte, par des échanges de consensus en HTTPS sur 2380. Un
        Patroni détient le verrou du leader, et une réplique éligible peut
        être promue si le primaire est perdu.
  \item \textbf{B -- Partition 2 + 1} : le côté majoritaire garde sa capacité
        de décision ; une promotion y reste possible sous conditions. Le
        membre isolé ne peut pas renouveler seul le verrou. S'il portait le
        primaire, Patroni tente d'arrêter ou de rétrograder son PostgreSQL.
  \item \textbf{C -- Deux membres perdus} : la majorité est indisponible et
        la coordination normale impossible. Aucune promotion coordonnée n'a
        lieu, et le dernier serveur passe en lecture seule. Rétablir le
        quorum exige un diagnostic.
  \item \textbf{Protection} : \code{watchdog.mode = off} dans les LXC ; aucun
        fencing matériel n'est configuré. Le quorum etcd ne garantit pas que
        toutes les transactions SQL sont répliquées.
\end{itemize}
\end{lecture}

\subsection{Bascule planifiée}

\schema{08-bascule-planifiee}{Bascule planifiée pour maintenance}

\begin{lecture}{Un opérateur choisit un candidat sain ; Patroni organise le changement de rôle du cluster.}
\begin{enumerate}
  \item \emph{Préparer} : vérifier que le quorum etcd est disponible, que le
        primaire est sain et que les répliques sont à jour, puis contrôler
        l'état et le retard du candidat.
  \item \emph{Demander la bascule} :
\begin{lstlisting}
patronictl -c /etc/patroni/config.yml switchover odoo-pg --candidate odoodb2
\end{lstlisting}
  \item \emph{Changer de primaire} : l'ancien leader se retire et le candidat
        obtient le rôle de primaire. Les clients peuvent perdre leur session.
\end{enumerate}
\textbf{Résultat :} l'ancien primaire odoodb1 devient réplique ; on vérifie la
réplication avant toute maintenance, et on intervient sur un nœud à la fois.
HAProxy suit \code{/primary} : le nouveau primaire odoodb2 répond 200 et
reçoit les nouvelles connexions ; odoodb3 le suit. On revient à l'état stable
d'un primaire et deux répliques.

\textbf{À contrôler :} un primaire, deux répliques, des connexions via chaque
HAProxy et des écritures applicatives. Une courte interruption des connexions
est possible. Le cron d'Odoo reste porté par odoo1, indépendamment du
primaire choisi par Patroni.\end{lecture}


% ===========================================================================
\section{Réseau et sécurité des flux}
% ===========================================================================

Chaque conteneur est relié au LAN en \textbf{IPVLAN L2}, sur la carte de son
hôte. Le réseau des hôtes rejette les adresses MAC inconnues : macvlan et veth
sur bridge n'obtenaient aucune réponse de la passerelle, alors qu'IPVLAN
réutilise la MAC de l'hôte. L'hôte garde sa carte, qui porte sa route par
défaut. En contrepartie, \textbf{un hôte ne joint pas son propre conteneur} :
les contrôles se font depuis le poste de contrôle, ou par \code{lxc-attach}.

\begin{table}[H]
\centering\small
\begin{tabularx}{\textwidth}{P{3.2cm} L}
\toprule
\rowcolor{gray!15}
\textbf{Port} & \textbf{Flux} \\
\midrule
22/TCP             & Poste de contrôle vers les hôtes et les conteneurs \\
8069/TCP           & Utilisateurs vers chaque Odoo (interface web) \\
8072/TCP           & Utilisateurs vers chaque Odoo (bus de notifications) \\
5000/TCP           & Odoo et sauvegarde vers les trois HAProxy \\
5432/TCP           & HAProxy vers PostgreSQL ; réplication et \code{pg_rewind} entre serveurs \\
8008/HTTP          & HAProxy vers l'API Patroni ; bascules \\
2379, 2380/HTTPS   & etcd : clients (Patroni) et échanges entre membres \\
7000/HTTP          & Page d'état HAProxy, sur \code{127.0.0.1} seulement \\
24007, 49152+/TCP  & Conteneurs Odoo et hôte de sauvegarde vers les serveurs GlusterFS \\
\bottomrule
\end{tabularx}
\caption{Ports et flux}
\end{table}

Aucun répartiteur HTTP n'est déployé : chaque utilisateur se connecte à
l'adresse d'une instance. Une configuration nginx est fournie à part
(\code{nginx/odoo-repartition.conf}), mais aucun playbook ne l'installe.

\schema{11-flux-reseau-et-tls}{Flux réseau, certificats et surfaces de gestion}

\begin{lecture}{Deux autorités de certification séparent les accès SQL de la coordination etcd.}
\begin{itemize}
  \item \textbf{Odoo} (réseau .36.0/24) : certificat client dans
        \code{/etc/odoo/pg-ssl}, \code{PGSSLMODE = verify-full}, autorité SQL
        \code{prosys-pg-ca}.
  \item \textbf{HAProxy} (.33.161 à .163) : frontal TCP sur 5000, relais
        PostgreSQL transparent, statistiques sur \code{127.0.0.1:7000}.
  \item \textbf{PostgreSQL} (5432) : TLS jusqu'au serveur SQL, mot de passe et
        certificat d'Odoo exigés, réplication des WAL en TLS.
  \item \textbf{Patroni} (HTTP 8008) : \code{GET /primary} pour les sondes de
        HAProxy ; les actions exigent une authentification ; l'API reste en
        HTTP, non chiffrée.
  \item \textbf{etcd} (HTTPS 2379 et 2380) : autorité distincte
        \code{prosys-etcd-ca}, certificats clients pour 2379, certificats de
        pairs pour 2380. Avec une autorité commune, le certificat client d'un
        conteneur Odoo ouvrirait etcd, donc la configuration de PostgreSQL.
  \item \textbf{Poste de contrôle Ansible} : SSH sur 22 vers les hôtes et les
        conteneurs ; il conserve la PKI et les secrets.
  \item \textbf{Identité du client SQL} : le primaire voit l'adresse du
        HAProxy, mais le certificat TLS continue d'identifier Odoo. Les règles
        \code{pg_hba.conf} autorisent les adresses des nœuds de données.
\end{itemize}
\end{lecture}

% ===========================================================================


% ===========================================================================
\section{Installation et déploiement d'Odoo}
% ===========================================================================

\subsection{Prérequis du poste de contrôle}

\begin{itemize}
  \item \textbf{Ansible 2.10.x} et les collections épinglées dans
        \code{requirements.yml} (\code{community.general 1.3.6},
        \code{community.postgresql 1.1.1}, \code{ansible.posix 1.1.1}).
        \code{ansible.posix 1.1.1} refuse de se charger avec un Ansible plus
        récent.
  \item \textbf{Accès SSH par clé} aux six hôtes ; le rôle \code{lxc_lan}
        injecte \code{~/.ssh/id_rsa.pub} dans les conteneurs.
  \item \textbf{Secrets} chiffrés dans \code{group_vars/all/vault.yml}.
  \item \textbf{Autorités de certification} dans \code{pki/}, sur le poste de
        contrôle et hors du dépôt. Le rôle \code{pg_ha} refuse d'en créer une
        nouvelle : il faut lancer les playbooks depuis le poste qui la détient.
  \item \textbf{Hors du dépôt} : le volume \code{sre} accepte les adresses des
        conteneurs Odoo, et les ports du tableau des flux sont ouverts.
\end{itemize}

\begin{table}[H]
\centering\small
\begin{tabularx}{\textwidth}{P{6.6cm} L}
\toprule
\rowcolor{gray!15}
\textbf{Secret} & \textbf{Usage} \\
\midrule
\code{vault_odoo_admin_password}           & Mot de passe maître d'Odoo (gestion des bases) \\
\code{vault_odoo_db_password}              & Rôle PostgreSQL \code{odoo}, partagé par la base et les Odoo \\
\code{vault_pg_replication_password}       & Rôle \code{replicator} (flux WAL) \\
\code{vault_pg_patroni_superuser_password} & Superutilisateur \code{postgres}, pour \code{pg_rewind} \\
\code{vault_pg_patroni_rest_password}      & Compte \code{patroni} de l'API REST \\
\bottomrule
\end{tabularx}
\caption{Secrets du vault}
\end{table}

\subsection{Ordre de déploiement}

L'ordre est contraignant :

\begin{lstlisting}
ansible-playbook playbooks/ansi2-snapchot.yaml                # 1. conteneurs LXC Debian 12
ansible-playbook playbooks/reseau-lxc.yaml --check --diff     # 2. voir les changements reseau
ansible-playbook playbooks/reseau-lxc.yaml                    #    IPVLAN + SSH (redemarre les LXC)
ansible-playbook playbooks/ansi_odoo.yaml --ask-vault-pass    # 3. base, puis Odoo
\end{lstlisting}

Le rôle de chaque playbook, ses étapes et ses précautions sont détaillés à la
section~\ref{sec:playbooks}. Le schéma suivant résume le parcours
d'installation d'Odoo que réalise \code{ansi_odoo.yaml}.

\schema{13-installation-odoo14-debian-lxc}{Installer Odoo 14 sur Debian dans les LXC}

\begin{lecture}{Parcours du projet : stockage et base prêts avant une initialisation unique et le démarrage des trois Odoo.}
\begin{enumerate}
  \item \emph{Debian~12 et LXC} : hôtes .36.21, .22 et .23 ; conteneurs
        Debian bookworm .36.161, .162 et .163 ; réseau IPVLAN L2.
  \item \emph{Filestore partagé} : client GlusterFS dans chaque conteneur,
        volume \code{192.168.36.8:/sre} monté sur \code{/var/lib/odoo},
        vérifié avant Odoo.
  \item \emph{PostgreSQL HA} : trois nœuds PostgreSQL~15 avec Patroni, etcd et
        HAProxy ; Odoo vise les trois adresses sur 5000, en TLS avec
        certificat client.
  \item \emph{Python et paquet Odoo} : CPython~3.8.20 compilé dans
        \code{/opt/python3.8} ; Odoo~14.0 par le dépôt APT nightly ; client
        PostgreSQL et wkhtmltopdf.
  \item \emph{Construire le venv} : \code{/opt/odoo14-venv}, dépendances pip
        dans le venv ; \code{odoo-deb.pth} rend visible le code fourni par le
        paquet \code{.deb}.
  \item \emph{Configurer les accès} : \code{/etc/odoo/odoo.conf}, certificats
        dans \code{/etc/odoo/pg-ssl}, fragment systemd \code{venv.conf}, montage
        du filestore exigé.
  \item \emph{Initialiser une fois} : UID 105 et GID 112 alignés, services
        Odoo arrêtés ; odoo1 seul initialise la base partagée \code{prod},
        et seulement si elle ne l'est pas déjà.
  \item \emph{Démarrer et vérifier} : systemd lance le Python du venv ; HTTP
        8069 et bus 8072 ; \code{workers=2} ; cron sur odoo1 seul ; imports
        d'Odoo et \code{safe_eval} validés.
\end{enumerate}
Le venv est local à chaque conteneur :
\code{/opt/python3.8/bin/python3.8 -m venv /opt/odoo14-venv}. Les dépendances
viennent du fichier figé s'il existe, sinon du \code{requirements.txt} amont
d'Odoo~14.
\end{lecture}

\subsection{Pile d'exécution d'Odoo}

Odoo~14 ne fonctionne pas sur le Python~3.11 de Debian~12 : son
\code{safe_eval} rejette le bytecode de Python~3.11 (opcode \code{RESUME}).
Aucune version de Debian ne fournit Python~3.8 ; il est donc compilé depuis
les sources et confiné dans \code{/opt}.

\begin{table}[H]
\centering\small
\begin{tabularx}{\textwidth}{P{3.2cm} L}
\toprule
\rowcolor{gray!15}
\textbf{Couche} & \textbf{Choix} \\
\midrule
Hôtes et conteneurs & Debian~12 bookworm, LXC \\
Interpréteur d'Odoo & CPython~3.8.20 dans \code{/opt/python3.8}, venv \code{/opt/odoo14-venv} \\
Odoo                & Paquet \code{.deb} nightly 14.0, lancé par le Python du venv \\
Base                & PostgreSQL~15 de Debian, sous Patroni \\
Coordination        & etcd, trois membres, quorum de deux \\
Routage SQL         & Un HAProxy par conteneur PostgreSQL, port 5000 \\
Filestore           & Client GlusterFS dans chaque conteneur Odoo, sur \code{/var/lib/odoo} \\
\bottomrule
\end{tabularx}
\caption{Pile logicielle}
\end{table}

\schema{14-pile-execution-odoo14-venv}{Pile d'exécution : LXC, venv et service Odoo}

\begin{lecture}{Le conteneur porte le système ; le venv fournit l'interpréteur Python et les bibliothèques de l'application.}
\begin{itemize}
  \item Emboîtement : hôte Debian~12 (exemple 192.168.36.21), conteneur LXC
        odoo1 (192.168.36.161), puis environnement virtuel Python.
  \item \textbf{systemd} lance \code{odoo.service}, complété par le fragment
        \code{venv.conf} :
\begin{lstlisting}
/opt/odoo14-venv/bin/python /usr/bin/odoo --config /etc/odoo/odoo.conf --logfile /var/log/odoo/odoo.log
\end{lstlisting}
  \item \textbf{Python système 3.11} (\code{/usr/bin/python3}) reste celui du
        système et d'Ansible.
  \item \textbf{Venv} : CPython~3.8.20 dans \code{/opt/python3.8}, venv
        \code{/opt/odoo14-venv}, dépendances pip dans
        \code{lib/python3.8/site-packages} ; \code{odoo-deb.pth} pointe vers le
        code du paquet.
  \item \textbf{Code d'Odoo~14} : fourni par le paquet \code{.deb}
        (\code{/usr/bin/odoo}, \code{/usr/lib/python3/dist-packages/odoo}, et
        les addons du paquet).
  \item \textbf{Configuration locale} : \code{/etc/odoo/odoo.conf} et
        certificats dans \code{/etc/odoo/pg-ssl}.
  \item \textbf{Données partagées} : GlusterFS monté dans le conteneur sur
        \code{/var/lib/odoo} (\code{data_dir}), avec \code{filestore/prod} et
        \code{sessions}.
  \item \textbf{Accès à la base} : Odoo, psycopg2 et libpq vers les trois
        HAProxy sur 5000, puis vers le primaire sur 5432. Le venv et le code
        sont locaux ; le \code{data_dir} est partagé.
\end{itemize}
\end{lecture}

% ===========================================================================
\section{Les playbooks en détail}
\label{sec:playbooks}
% ===========================================================================

Chaque playbook ne fait qu'orchestrer : la logique vit dans des rôles, rangés
par playbook sous \code{roles/<playbook>/<rôle>}. Le \code{roles_path} de
\code{ansible.cfg} les retrouve tous par leur seul nom. Les playbooks se lancent
depuis la racine du dépôt, avec l'inventaire unique \code{inventory.yaml}.

\begin{table}[H]
\centering\small
\begin{tabularx}{\textwidth}{P{4.6cm} P{3.1cm} L}
\toprule
\rowcolor{gray!15}
\textbf{Playbook} & \textbf{Cibles} & \textbf{Rôle et moment d'utilisation} \\
\midrule
\code{ansi2-snapchot.yaml}      & les six hôtes                  & Créer les conteneurs LXC ; une fois, avant tout le reste \\
\code{reseau-lxc.yaml}          & les six hôtes                  & Réseau IPVLAN et SSH des conteneurs ; après la création, ou pour réparer le réseau \\
\code{ansi_odoo.yaml}           & hôtes et conteneurs            & Déploiement complet ; rejouable pour remettre en conformité \\
\code{filestore-conteneur.yaml} & conteneurs et hôtes Odoo       & Montage GlusterFS par les conteneurs ; joué aussi par \code{ansi_odoo.yaml} et \code{ansi_syncro.yaml} \\
\code{ansi_postgresql_ha.yaml}  & les trois serveurs PostgreSQL  & Reprise du cluster par Patroni, puis remise en conformité \\
\code{ansi_syncro.yaml}         & hôtes, conteneurs et base      & Migration des pièces jointes entre la base et le filestore \\
\code{nettoyage_odoo.yaml}      & conteneurs Odoo                & Réparation APT après un lancement interrompu \\
\code{nettoyage.yaml}           & les six hôtes                  & Destruction des conteneurs et désinstallation de LXC \\
\bottomrule
\end{tabularx}
\caption{Vue d'ensemble des playbooks}
\end{table}

Le playbook \code{ansi_sauvegarde.yaml}, qui déployait la sauvegarde, n'est plus
présent dans le dépôt (section~\ref{sec:sauvegarde}).

% ---------------------------------------------------------------------------
\subsection{\pb{ansi2-snapchot.yaml} -- créer les conteneurs}

\paragraph{Objectif.} Installer LXC sur chaque hôte et y créer le conteneur
Debian~12 déclaré dans \code{host_vars/} : \code{devodoo14} sur les hôtes Odoo,
\code{devpg15} sur les hôtes PostgreSQL. La liste des conteneurs vient de
\code{common/group_vars/local_dev.yml}, qui la dérive de
\code{lxc_container_name} hôte par hôte.

\paragraph{Déroulement} (rôle \code{install}, sur chaque hôte).
\begin{enumerate}
  \item Détecter le système et installer les paquets LXC. Une branche
        Rocky/RHEL (EPEL, SELinux permissif) est conservée pour l'ancien lab.
  \item Démarrer le réseau LXC (\code{lxcbr0}) et le service LXC, puis
        détruire les conteneurs de diagnostic laissés par un essai précédent.
  \item Mesurer l'espace libre sous \code{/var/lib/lxc}. Si la partition est
        trop petite, déporter le stockage LXC sous \code{/srv} : copie du
        contenu, montage et inscription dans \code{/etc/fstab}. Le rôle exige
        ensuite assez d'espace pour le système du conteneur.
  \item Supprimer les conteneurs obsolètes listés dans \code{stale_containers},
        si cette variable est définie.
  \item Créer et démarrer le conteneur Debian~12 \emph{bookworm} depuis le
        modèle \code{download}. En cas d'échec, le rôle rejoue
        \code{lxc-create} sans \code{--quiet} et affiche l'état du disque, du
        réseau et le journal de création.
  \item Activer le démarrage automatique du conteneur.
\end{enumerate}

\paragraph{Lancement.}
\begin{lstlisting}
ansible-playbook playbooks/ansi2-snapchot.yaml
\end{lstlisting}

\paragraph{Précautions.} Le conteneur créé n'a pas encore son réseau
définitif : il faut enchaîner avec \code{reseau-lxc.yaml}. Renseigner
\code{stale_containers} détruit des conteneurs : sauvegarder avant.

% ---------------------------------------------------------------------------
\subsection{\pb{reseau-lxc.yaml} -- relier les conteneurs au réseau}

\paragraph{Objectif.} Exposer chaque conteneur directement sur le LAN en
IPVLAN L2, puis amorcer son accès SSH pour qu'Ansible puisse le gérer. C'est le
seul playbook qui joue le rôle \code{lxc_lan} : \code{ansi_odoo.yaml} suppose
qu'il a déjà été passé.

\paragraph{Déroulement} (rôle \code{lxc_lan}, sur chaque hôte).
\begin{enumerate}
  \item Contrôler la description : mode réseau et adressage présents dans
        \code{host_vars/}, passerelle dans le sous-réseau du conteneur, carte
        réseau présente. Le rôle refuse de céder au conteneur la carte qui
        porte la route par défaut de l'hôte.
  \item Sauvegarder, une seule fois, la configuration LXC d'origine, puis
        retirer les réglages des anciens modes (macvlan, adresse MAC figée).
  \item Écrire l'interface du conteneur en \code{ipvlan} sur la carte de
        l'hôte, et son adresse IP statique dans le conteneur
        (\code{systemd-networkd}, sans DHCP). La sortie NAT de secours par
        \code{lxcbr0} est retirée, car \code{lxc_egress_nat=false}.
  \item Arrêter puis redémarrer le conteneur, attendre l'état \code{RUNNING}
        et la configuration d'\code{eth0}. En cas d'échec, le rôle affiche la
        configuration et le journal LXC.
  \item Installer \code{openssh-server}, \code{python3} et \code{python3-apt},
        déposer la clé publique du poste de contrôle, puis attendre que le
        port 22 du conteneur réponde, vu du poste de contrôle.
\end{enumerate}

\paragraph{Lancement.}
\begin{lstlisting}
ansible-playbook playbooks/reseau-lxc.yaml --check --diff        # voir ce qui changerait
ansible-playbook playbooks/reseau-lxc.yaml
ansible-playbook playbooks/reseau-lxc.yaml --limit 192.168.36.21  # un seul hote
\end{lstlisting}

\paragraph{Précautions.} Le playbook \textbf{redémarre les conteneurs}. En
IPVLAN, un hôte ne joint pas son propre conteneur : les tests se font depuis
le poste de contrôle. Le poste doit avoir la version à jour du rôle : une
version ancienne réécrirait le mode macvlan et casserait un conteneur déjà
passé en IPVLAN. La procédure manuelle équivalente est décrite dans
\code{docs/BOOTSTRAP-CONTENEUR-MANUEL.md}.

% ---------------------------------------------------------------------------
\subsection{\pb{ansi\_odoo.yaml} -- déploiement complet}

\paragraph{Objectif.} Déployer la base, puis les trois instances Odoo, dans un
ordre contraignant. Rejoué sur une installation existante, il remet chaque
élément en conformité.

\paragraph{Déroulement.}

\begin{table}[H]
\centering\small
\begin{tabularx}{\textwidth}{P{1.7cm} P{3.1cm} L}
\toprule
\rowcolor{gray!15}
\textbf{Play} & \textbf{Cible} & \textbf{Effet} \\
\midrule
0/5         & poste de contrôle & Contrôle d'environnement : Ansible entre 2.10 et 2.11, collections de \code{requirements.yml} présentes. Tag \code{never} : à demander par \code{--tags all,preflight} \\
1/5         & les six hôtes     & \code{chrony} remet l'horloge à l'heure : une horloge fausse fait échouer \code{apt} dans les conteneurs \\
1bis/5      & conteneurs, puis hôtes Odoo & Import de \code{filestore-conteneur.yaml} : le volume doit être monté avant l'installation d'Odoo \\
2/5, 2bis/5 & \code{db_containers}, puis \code{db_replicas} & \code{pg_server} : primaire, puis réserves en streaming. Plays sautées quand \code{pg_ha_enabled: true} \\
2ter/5      & \code{db_cluster} & \code{pg_ha} : Patroni, etcd et HAProxy sur les trois serveurs ensemble ; un échec arrête les trois. Sur un cluster déjà repris, ne fait que vérifier \\
3/5         & \code{odoo_containers} & \code{odoo14}, puis \code{odoo_uid} : Odoo installé et configuré, service laissé arrêté \\
4/5         & \code{odoo_primary} & \code{odoo_dbinit} : \code{odoo -i base --without-demo=all} sur la base partagée, par odoo1 seul, avec son certificat client \\
5/5         & \code{odoo_containers} & Démarrage et activation d'Odoo, attente du port 8069 (120~s au plus) \\
\bottomrule
\end{tabularx}
\caption{Plays du playbook \texttt{ansi\_odoo.yaml}}
\end{table}

Le play 3/5 fait l'essentiel du travail dans chaque conteneur Odoo :
\begin{itemize}
  \item sources APT cohérentes (rôle \code{base_debian}) et CPython~3.8.20
        compilé dans \code{/opt/python3.8} (rôle \code{python38}) ;
  \item dépôt APT nightly d'Odoo~14 et installation du paquet \code{.deb} ;
  \item environnement virtuel \code{/opt/odoo14-venv} sur Python~3.8. Les
        dépendances viennent d'un fichier de versions figées s'il existe sur
        le poste de contrôle, sinon du \code{requirements.txt} amont ;
  \item \code{odoo-deb.pth} rend le code du paquet visible depuis le venv ; une
        enveloppe \code{odoo} et un fragment systemd font lancer le service
        par le Python du venv ;
  \item refus de configurer Odoo si \code{/var/lib/odoo} n'est pas le volume
        GlusterFS, et exigence de ce montage au démarrage ;
  \item vérification qu'Odoo s'importe sous Python~3.8 et que \code{safe_eval}
        accepte ses expressions ;
  \item dépôt du certificat client du nœud dans \code{/etc/odoo/pg-ssl}, puis
        de \code{/etc/odoo/odoo.conf} ;
  \item contrôle que la base partagée écoute et que ce nœud est autorisé à s'y
        connecter ; Odoo reste arrêté jusqu'à l'initialisation.
\end{itemize}
Le rôle \code{odoo_uid} aligne ensuite l'UID et le GID du compte \code{odoo}
(105/112) sur les trois conteneurs. Trois \code{odoo -i base} concurrents sur
la même base corrompraient le registre : c'est pour cela que l'initialisation
(4/5) est confiée à odoo1 seul et précède le démarrage commun (5/5).

\paragraph{Lancement.}
\begin{lstlisting}
ansible-playbook playbooks/ansi_odoo.yaml --ask-vault-pass
ansible-playbook playbooks/ansi_odoo.yaml --ask-vault-pass --tags all,preflight  # avec le controle du poste
ansible-playbook playbooks/ansi_odoo.yaml --ask-vault-pass --tags horloge        # horloge seule
ansible-playbook playbooks/ansi_odoo.yaml --ask-vault-pass --tags filestore      # montage seul
\end{lstlisting}

\paragraph{Précautions.} Les conteneurs doivent exister et avoir leur réseau
(\code{ansi2-snapchot.yaml}, puis \code{reseau-lxc.yaml}). Le vault doit être
renseigné, et le volume GlusterFS doit accepter les adresses des conteneurs
Odoo. Le commentaire d'en-tête du playbook cite encore les adresses de
l'ancien lab (192.168.96.x) : c'est l'inventaire qui fait foi.

% ---------------------------------------------------------------------------

% ---------------------------------------------------------------------------
\subsection{\pb{ansi\_postgresql\_ha.yaml} -- haute disponibilité de la base}

\paragraph{Objectif.} Installer etcd, Patroni et HAProxy dans les trois
conteneurs PostgreSQL. Au premier passage, il reprend le cluster installé par
\code{pg_server}, sans perte ni recopie. Aux passages suivants, il remet en
conformité, sans coupure.

\paragraph{Déroulement} (rôle \code{pg_ha}, sur les trois serveurs ensemble).
\begin{enumerate}
  \item \emph{Contrôles} : trois serveurs, secrets présents, autorité de
        certification existante. Le rôle refuse d'en créer une nouvelle, car
        Odoo rejetterait alors tous les serveurs.
  \item \emph{Paquets} : \code{patroni}, \code{etcd-server} et \code{haproxy},
        installés sans démarrer leurs services.
  \item \emph{Constat et décision} : l'état de chaque serveur est lu sans rien
        arrêter (tableau~\ref{tab:decision}).
  \item \emph{Certificats et etcd} : certificats PostgreSQL et autorité etcd
        distincte ; etcd démarré sur les trois serveurs en même temps, puis
        vérifié sain avec ses trois membres.
  \item \emph{Patroni et reprise} : configuration de Patroni, puis reprise du
        primaire, confirmée par etcd, puis des réserves. Le PostgreSQL de
        Debian est retiré : unité masquée, \code{start.conf} désactivé,
        anciens fichiers de \code{pg_server} renommés en \code{.avant-patroni}.
  \item \emph{Relance} de Patroni, en pause et un serveur à la fois, seulement
        s'il tourne avec une ancienne liste etcd.
  \item \emph{Comptes et configuration} : mot de passe de \code{postgres} (pour
        \code{pg_rewind}), rôles \code{replicator} et \code{odoo}, base
        \code{prod} ; configuration dynamique d'etcd remise en conformité avec
        \code{pg_ha_dcs}, en conservant une pause posée à la main.
  \item \emph{HAProxy, puis vérification} : chaque Patroni joint les trois
        membres etcd ; un primaire et deux réserves avec moins de 16~Mo de
        retard ; une vraie connexion TLS à travers chaque HAProxy aboutit au
        primaire.
\end{enumerate}

\begin{table}[H]
\centering\small
\begin{tabularx}{\textwidth}{P{2cm} P{5.4cm} L}
\toprule
\rowcolor{gray!15}
\textbf{Décision} & \textbf{Quand} & \textbf{Effet} \\
\midrule
\code{demarrer} & serveur déjà repris par Patroni & Patroni est démarré s'il ne tourne pas \\
\code{primaire} & premier passage, odoodb1, qui n'est pas en réserve & Debian arrête PostgreSQL ; Patroni le redémarre et adopte ses données \\
\code{reserve}  & même cluster que le primaire, en réserve & repris en réserve, sans recopie \\
\code{recopie}  & cluster étranger sans aucune base & vidé puis recopié depuis le primaire \\
\code{refus}    & tout autre cas & arrêt du playbook, avant toute modification \\
\bottomrule
\end{tabularx}
\caption{Décision de reprise, par serveur}
\label{tab:decision}
\end{table}

\paragraph{Lancement.}
\begin{lstlisting}
ansible-playbook playbooks/ansi_postgresql_ha.yaml --ask-vault-pass
\end{lstlisting}

\paragraph{Précautions.} Jamais de \code{--limit} : etcd et la reprise exigent
les trois serveurs ensemble. Lancer depuis le poste qui détient \code{pki/} et
le vault, avec les deux secrets Patroni. Ouvrir 2379, 2380, 8008 et 5432 entre
les serveurs, et 5000 depuis le réseau Odoo. Au premier passage, Odoo perd la
base quelques secondes. Il faut ensuite passer sans attendre
\code{pg_ha_enabled: true}, puis rejouer \code{ansi_odoo.yaml} : sinon, après
une bascule, Odoo resterait sur un serveur en lecture seule.

% ---------------------------------------------------------------------------

% ---------------------------------------------------------------------------
\subsection{\pb{nettoyage\_odoo.yaml} -- réparer un conteneur Odoo}

\paragraph{Objectif.} Remettre à plat les conteneurs Odoo après un lancement
interrompu, sans les recréer.

\paragraph{Déroulement par défaut, non destructif} (rôle \code{odoo_cleanup}) :
arrêt d'Odoo ; libération des verrous \code{dpkg}/\code{apt} résiduels ;
reprise d'une transaction \code{dpkg} interrompue ; réparation des paquets
laissés incohérents ; réalignement des versions par l'outil
\code{apt-realign} ; réparation des dépendances ; vidage du cache ; résumé de
l'état du conteneur.

\paragraph{Option \texttt{full\_reset}.} Avec \code{-e full_reset=true}, le rôle
purge en plus Odoo et PostgreSQL, puis supprime leurs répertoires
(\code{/var/lib/odoo}, \code{/etc/odoo}, \code{/var/log/odoo},
\code{/var/lib/postgresql}, \code{/etc/postgresql}) et le dépôt nightly.

\begin{tcolorbox}[colback=red!4, colframe=red!60!black, boxrule=0.6pt,
                  fonttitle=\bfseries\small, title={Risque : \texttt{full\_reset} et filestore partagé}]
\small
Dans l'architecture actuelle, \code{/var/lib/odoo} est le point de montage du
volume GlusterFS commun. Supprimer ce répertoire effacerait le contenu du
filestore partagé par les trois instances. Le garde-fou
\code{destroy_shared_db} ne protège que les hôtes du groupe
\code{db_containers}, que ce playbook ne vise pas. Ne pas utiliser
\code{full_reset} en l'état sans avoir d'abord démonté le volume et vérifié une
sauvegarde.
\end{tcolorbox}

\paragraph{Lancement.}
\begin{lstlisting}
ansible-playbook playbooks/nettoyage_odoo.yaml --ask-vault-pass                     # reparation APT
ansible-playbook playbooks/nettoyage_odoo.yaml --ask-vault-pass -e full_reset=true  # DESTRUCTIF
\end{lstlisting}

% ---------------------------------------------------------------------------
\subsection{\pb{nettoyage.yaml} -- tout désinstaller}

\paragraph{Objectif.} Repartir de zéro au niveau des hôtes.

\paragraph{Déroulement} (rôle \code{cleanup}, sur les six hôtes) : destruction
des conteneurs LXC ; arrêt et désactivation des services LXC ; suppression des
paquets, des répertoires et des fichiers créés pour LXC ; en option
(\code{-e remove_epel=true}), suppression du dépôt EPEL ; sur Rocky/RHEL, retour
de SELinux en mode \code{enforcing}.

\paragraph{Lancement.}
\begin{lstlisting}
ansible-playbook playbooks/nettoyage.yaml
\end{lstlisting}

\paragraph{Précautions.} Le groupe visé comprend aussi les trois hôtes
PostgreSQL : le playbook détruit les conteneurs \code{devpg15}, donc les
données de la base. Sauvegarder d'abord. Pour repartir de zéro, enchaîner
\code{nettoyage.yaml}, \code{ansi2-snapchot.yaml}, \code{reseau-lxc.yaml}, puis
\code{ansi_odoo.yaml} et la séquence de passage à Patroni.

% ===========================================================================
\section{Vérification et exploitation}
% ===========================================================================

\subsection{Vérification après déploiement}

Depuis le poste de contrôle, puisqu'un hôte ne joint pas son propre
conteneur :

\begin{lstlisting}
# Odoo actif sur les trois conteneurs
ansible odoo_containers -m shell -a 'systemctl is-active odoo'

# /var/lib/odoo est le volume GlusterFS (attendu : fuse.glusterfs 192.168.36.8:/sre)
ansible odoo_containers -m shell -a 'findmnt -n -o FSTYPE,SOURCE /var/lib/odoo'

# un Leader et deux Replica en streaming, retard proche de 0
ansible odoodb1 -m shell -a 'patronictl -c /etc/patroni/config.yml list'

# interface web des trois instances (attendu : 200)
for i in 161 162 163; do curl -sI -o /dev/null -w "$i -> %{http_code}\n" \
  http://192.168.36.$i:8069/web/login; done
\end{lstlisting}

Le compte créé par \code{odoo -i base} est \code{admin} / \code{admin} ; son
mot de passe est à changer dès la première connexion. Le mot de passe du vault
\code{vault_odoo_admin_password} est le mot de passe maître de gestion des
bases, et le gestionnaire de bases web est désactivé (\code{list_db = false}).

\subsection{Exploitation courante}

Sur un conteneur PostgreSQL :

\begin{lstlisting}
patronictl -c /etc/patroni/config.yml list                     # etat du cluster
patronictl -c /etc/patroni/config.yml switchover odoo-pg       # bascule planifiee
patronictl -c /etc/patroni/config.yml restart odoo-pg          # redemarrer PostgreSQL
patronictl -c /etc/patroni/config.yml reinit odoo-pg odoodb3   # recopier une reserve
journalctl -u patroni -f                                       # journal de Patroni
ssh -L 7000:127.0.0.1:7000 root@192.168.33.161                 # page d'etat HAProxy
\end{lstlisting}

\begin{itemize}
  \item C'est Patroni qui démarre et arrête PostgreSQL : on n'utilise plus
        \code{systemctl restart postgresql}, \code{pg_ctlcluster} ni
        \code{pg_ctl}, et l'unité Debian est masquée.
  \item Un paramètre commun se change dans \code{pg_ha_dcs}, puis on rejoue le
        playbook. Un \code{patronictl edit-config} serait écrasé au passage
        suivant.
  \item \code{/etc/patroni/config.yml} ne se corrige pas à la main : on rejoue
        le playbook, qui relance Patroni en pause si la liste etcd a changé.
\end{itemize}

% ===========================================================================
\section{Évolution de l'architecture}
\label{sec:evolution}
% ===========================================================================

Les schémas 15 à 24 ont été produits avant l'édition du 3~novembre~2026. Ils
décrivent des états antérieurs du dépôt : haute disponibilité non activée,
GlusterFS monté sur l'hôte puis passé au conteneur par un bind. Ils sont
conservés parce qu'ils expliquent les choix actuels. Chaque lecture précise ce
qui a changé depuis ; en cas de divergence, les schémas 1 à 14 font foi.

% ---------------------------------------------------------------------------
\subsection{Réseau des conteneurs LXC }

Ces trois vues décrivent le raccordement IPVLAN, toujours en vigueur, et les
chemins réseau d'un utilisateur et d'Odoo.

\schema{schemas-lxc/01-raccordement-ipvlan}{Raccordement IPVLAN : un conteneur, une IP dédiée}

\begin{lecture}{IPVLAN L2 partage la carte réseau et l'adresse MAC de l'hôte.}
\begin{itemize}
  \item L'hôte LXC 192.168.36.21 (Debian~12) garde sa carte \code{enp5s0} et
        son adresse propre, 192.168.36.21.
  \item Le conteneur \code{devodoo14} (odoo1) reçoit une adresse statique :
        \code{eth0} en 192.168.36.161/24, passerelle 192.168.36.254, DNS
        172.24.10.17 et 172.24.10.18, configurés dans le conteneur.
  \item L'adresse de l'hôte et l'interface IPVLAN passent par la même carte.
        Sur le LAN applications (192.168.36.0/24), le réseau voit donc deux
        adresses IP derrière une seule adresse MAC : les deux adresses doivent
        y être autorisées.
  \item Le raccordement est direct : ni bridge Linux, ni NAT sur l'hôte.
\end{itemize}

\end{lecture}

\schema{schemas-lxc/02-acces-utilisateur-odoo}{Accès utilisateur : du navigateur à Odoo}

\begin{lecture}{La destination 192.168.36.161:8069 conduit directement au service dans le conteneur.}
\begin{enumerate}
  \item \emph{Adresse ciblée} : depuis un poste du même LAN, ou disposant
        d'une route, le navigateur ouvre \code{http://192.168.36.161:8069}.
  \item \emph{Distribution par Linux} : le paquet arrive sur la carte
        \code{enp5s0} de l'hôte 192.168.36.21. Linux le distribue selon son
        adresse de destination : IPVLAN remet ce qui est destiné à .36.161 à
        l'interface \code{eth0} du conteneur.
  \item \emph{Service Odoo} : Odoo~14 répond en HTTP sur TCP/8069 ; le bus de
        notifications (longpolling) utilise TCP/8072. La réponse HTTP repart
        vers le navigateur.
\end{enumerate}

\end{lecture}

\schema{schemas-lxc/03-flux-odoo-postgresql}{Flux Odoo vers PostgreSQL : deux sous-réseaux}

\begin{lecture}{Deux sous-réseaux, une connexion SQL vers le primaire PostgreSQL.}
\begin{itemize}
  \item Côté applications : odoo1 (\code{devodoo14}, \code{eth0} en
        192.168.36.161/24, hôte .36.21 par \code{enp5s0}) vise
        192.168.33.161:5432.
  \item Côté bases de données : odoodb1 (\code{devpg15}, \code{eth0} en
        192.168.33.161/24, hôte .33.21 par \code{enp1s0}) porte le
        PostgreSQL~15 primaire, à l'écoute sur TCP/5432.
  \item Entre les deux, un ou plusieurs équipements L3 externes assurent le
        routage. Les passerelles .36.254 et .33.254 sont déclarées dans
        \code{host_vars/} ; leur implantation physique est hors du dépôt.
  \item Conditions de fonctionnement : routage aller et retour, autorisation
        de TCP/5432, écoute de PostgreSQL sur .33.161 et autorisation des
        clients Odoo dans \code{pg_hba.conf}. Le NAT est désactivé sur les
        hôtes.
\end{itemize}

\end{lecture}

% ---------------------------------------------------------------------------
\subsection{Filestore GlusterFS }

Ces trois vues décrivent le volume partagé et l'ancien modèle de montage, dans
lequel l'hôte montait le volume puis le passait au conteneur.

\schema{schemas-glusterfs/01-volume-partage-trois-odoo}{Volume GlusterFS partagé entre les trois Odoo}

\begin{lecture}{Un stockage logique unique pour les pièces jointes des trois instances.}
\begin{itemize}
  \item \textbf{Périmètre} : le volume GlusterFS existe déjà, hors du projet.
        Le projet configure ses clients et son partage avec les conteneurs.
  \item \textbf{Volume externe} \code{sre} : la description du volume est
        lue auprès de \code{192.168.36.8:/sre} ; les briques sont hébergées
        hors du projet.
  \item Après lecture de la description, chaque client contacte directement
        les serveurs de briques du volume.
  \item \textbf{Stockage partagé} : une pièce jointe écrite par odoo1 devient
        accessible à odoo2 et odoo3.
  \item Le groupe \code{gluster_bricks} désigne les hôtes clients ; les hôtes
        PostgreSQL ne reçoivent pas ce montage.
\end{itemize}

\end{lecture}

\schema{schemas-glusterfs/02-montage-hote-et-bind-lxc}{Ancien modèle : montage sur l'hôte et bind LXC}

\begin{lecture}{Le volume monté sur l'hôte était passé au conteneur par un bind.}
\begin{itemize}
  \item Le client GlusterFS de l'hôte 192.168.36.21 montait
        \code{192.168.36.8:/sre} sur \code{/mnt/odoo}, de façon persistante
        par \code{/etc/fstab}.
  \item Un bind LXC présentait ce répertoire au conteneur odoo1
        (192.168.36.161) sur \code{/var/lib/odoo}. Ansible écrivait dans
        \code{/var/lib/lxc/devodoo14/config} la directive
        \code{lxc.mount.entry = /mnt/odoo var/lib/odoo none bind,create=dir 0 0}.
  \item \code{/var/lib/odoo} est le \code{data_dir} d'Odoo : il contient
        \code{filestore/prod/} (documents et images) et \code{sessions/}
        (sessions partagées).
  \item Garde-fous de l'époque : montage après le réseau
        (\code{defaults,_netdev}), contrôle du type \code{fuse.glusterfs} par
        \code{findmnt}, redémarrage du conteneur si la directive changeait.
\end{itemize}
\textbf{Pourquoi ce modèle a été abandonné :} un hôte dont le montage échouait
au démarrage passait au conteneur un \code{/mnt/odoo} vide. Le bind faisait de
\code{/var/lib/odoo} un point de montage, la garde systemd laissait Odoo
démarrer, et Odoo écrivait sur le disque local .
Désormais, le conteneur monte lui-même le volume, et Odoo ne démarre que s'il
est présent.
\end{lecture}

\schema{schemas-glusterfs/03-glusterfs-et-sauvegarde}{GlusterFS et sauvegarde}

\begin{lecture}{La sauvegarde réunit un dump de la base et une archive du volume, sur un seul hôte.}
\begin{itemize}
  \item L'hôte de sauvegarde 192.168.36.21 lit le volume monté sur
        \code{/mnt/odoo}.
  \item Le script enchaîne trois étapes : vérifier le montage GlusterFS, faire
        un \code{pg_dump} de la base, puis archiver les fichiers.
        \code{./sessions} et \code{./_sessions} sont exclus ; le reste de
        \code{/mnt/odoo} est archivé.
  \item Base : primaire 192.168.33.161:5432, base partagée \code{prod} par
        défaut, lue par \code{pg_dump} depuis 192.168.36.21.
  \item Les fichiers sont écrits dans \code{/srv/sauvegardes}, sur le disque
        local de .36.21, séparé du volume : \code{prod-<date>.dump},
        \code{filestore-<date>.tgz} et \code{manifeste-<date>.txt}.
  \item Précondition : le volume doit être déjà monté ; le script contrôle sa
        présence. Le schéma cite le playbook \code{ansi_sauvegarde.yaml}, depuis
        supprimé du dépôt.
\end{itemize}

\end{lecture}

% ---------------------------------------------------------------------------
\subsection{Architecture et HA PostgreSQL }

Ces quatre vues datent d'avant l'activation de Patroni : la base était servie
par \code{pg_server}, avec une bascule manuelle, et la haute disponibilité
n'était encore qu'une option.

\schema{schemas-architecture-ha/01-architecture-actuelle}{Vue de déploiement et flux de production (avant Patroni)}

\begin{lecture}{Six hôtes LXC ; vue logique des réseaux, des services et des dépendances, flux par flux.}
\begin{itemize}
  \item Acteurs : les utilisateurs Odoo, qui accèdent directement aux trois
        conteneurs, et le contrôleur Ansible, en SSH sur 22 vers les six
        hôtes et les six conteneurs.
  \item Zones : application (192.168.36.0/24, passerelle .36.254) et données
        (192.168.33.0/24, passerelle .33.254), en IPVLAN L2 sans NAT sur les
        hôtes.
  \item \textbf{F01} HTTP 8069 et longpolling 8072 ; \textbf{F02} SQL-TLS vers
        le primaire sur 5432, avec certificat client, à travers le routage L3
        et son SNAT vers .33.254 ; \textbf{F03} clients GlusterFS sur les hôtes
        vers \code{/mnt/odoo}, bind vers \code{/var/lib/odoo}, TCP/24007 et
        ports des briques ; \textbf{F04} streaming WAL vers deux répliques
        (slots \code{odoodb2} et \code{odoodb3}) ; \textbf{F05a} dump du
        primaire et \textbf{F05b} archive de \code{/mnt/odoo} sur .36.21 ;
        \textbf{F06} administration SSH.
  \item Disponibilité retenue à cette date : \code{pg_ha_enabled=false},
        primaire SQL et cron odoo1 uniques, bascule manuelle, réplication
        asynchrone.
  \item Option HA hors du chemin SQL : Patroni, etcd et HAProxy sur 5000 ;
        rôle présent, activation et validation à réaliser.
\end{itemize}

\end{lecture}

\schema{schemas-architecture-ha/02-replication-et-bascule-manuelle}{Réplication et bascule manuelle (avant Patroni)}

\begin{lecture}{Un primaire, deux répliques en streaming ; les clients n'écrivent que sur odoodb1.}
\begin{itemize}
  \item Mode déclaré : \code{pg_ha_enabled = false}, rôle \code{pg_server},
        aucune promotion automatique.
  \item odoo1, odoo2 et odoo3 se connectent tous à 192.168.33.161:5432 :
        odoodb1, sur l'hôte .33.21, reçoit toutes les lectures et écritures
        SQL.
  \item odoodb2 (.33.162, hôte .33.22) et odoodb3 (.33.163, hôte .33.23) sont
        des répliques en lecture seule, alimentées par WAL asynchrones à
        travers les slots \code{odoodb2} et \code{odoodb3}.
  \item Les réserves vivent sur deux hôtes distincts et conservent des copies
        de la base. Les dernières transactions peuvent manquer sur une
        réplique au moment d'une panne.
  \item En cas de panne du primaire : intervention manuelle, c'est-à-dire
        promouvoir une réserve puis changer \code{odoo_db_host}. Sans
        intervention, la panne d'odoodb1 interrompt les accès SQL d'Odoo,
        malgré les deux répliques.
\end{itemize}

\end{lecture}

\schema{schemas-architecture-ha/03-option-ha-patroni-etcd-haproxy}{Option HA Patroni, etcd et HAProxy (telle que prévue)}

\begin{lecture}{Patroni, etcd et trois points d'entrée SQL ; le placement initial n'est qu'un exemple.}
\begin{itemize}
  \item À cette date : option non activée (\code{pg_ha_enabled = false}) ;
        code présent, tests et déploiement à réaliser.
  \item Clients Odoo par libpq : \code{192.168.33.161,192.168.33.162,192.168.33.163}
        sur 5000, essayés dans l'ordre, avec
        \code{PGTARGETSESSIONATTRS=read-write}.
  \item etcd : trois membres, quorum requis de 2 sur 3.
  \item HAProxy interroge Patroni (\code{GET /primary}, 8008) ; Patroni gère
        PostgreSQL et le verrou du leader dans etcd.
  \item TLS SQL de bout en bout, streaming WAL asynchrone, aucun équilibrage
        des lectures vers les répliques, aucune adresse IP virtuelle.
  \item odoodb1 primaire est le placement initial ; Patroni peut désigner un
        autre nœud.
\end{itemize}

\end{lecture}

\schema{schemas-architecture-ha/04-bascule-ha-et-quorum}{Bascule HA et limites de disponibilité}

\begin{lecture}{Panne du primaire : élection, puis reconnexion. Scénario illustratif, non issu d'un test.}
\begin{itemize}
  \item odoodb1 (192.168.33.161) devient indisponible : PostgreSQL, Patroni,
        etcd et son HAProxy local sont perdus.
  \item odoodb2 (192.168.33.162) est élu primaire : promotion par Patroni sous
        conditions, etcd2 et etcd3 formant la majorité.
  \item odoodb3 (192.168.33.163) reste réplique et suit le nouveau primaire en
        streaming ; son HAProxy et son etcd restent disponibles.
  \item Les clients Odoo réessaient sur une adresse disponible ; les HAProxy
        .33.162 et .33.163 ciblent le primaire courant et relaient le SQL vers
        .33.162:5432.
  \item Quorum de 2 sur 3 : deux membres doivent communiquer.
  \item Limites : pertes possibles en asynchrone, pas de fencing matériel,
        durée à mesurer, certificats de sauvegarde en HA à compléter.
\end{itemize}

\end{lecture}

% ===========================================================================
\section{Limites connues}
% ===========================================================================

\begin{itemize}
  \item \textbf{Pas de chien de garde matériel} : un conteneur LXC n'a pas de
        \code{/dev/watchdog}. Un Patroni bloqué sans planter laisserait
        l'ancien primaire accepter des écritures pendant la promotion d'une
        réserve.
  \item \textbf{Réplication asynchrone} : lors d'une bascule non planifiée,
        des transactions récentes peuvent être perdues. Les transactions en
        cours ne sont pas transférées, et les clients doivent se reconnecter.
  \item \textbf{Quorum etcd de 2 sur 3} : la perte d'un serveur est tolérée.
        Si deux serveurs sont perdus, le dernier passe en lecture seule et Odoo
        s'arrête, par conception.
  \item \textbf{La haute disponibilité ne concerne que la base} : le cron
        d'Odoo reste sur odoo1, chaque accès web vise une instance, et le
        volume GlusterFS comme l'infrastructure physique sont hors du dépôt.
  \item \textbf{Groupes d'inventaire} : après une bascule, \code{db_containers}
        et \code{db_replicas} ne décrivent plus que le placement d'origine.
        \code{ansi_syncro.yaml} et \code{odoo_cleanup} traitent
        \code{db_containers} comme « la » base : il faut les vérifier avant de
        les lancer.
  \item \textbf{Fins de vie} : Odoo~14 n'est plus maintenu, et Python~3.8 ne
        reçoit plus de correctifs depuis octobre~2024. L'interpréteur est
        confiné à \code{/opt} et ne sert qu'à Odoo.
  \item \textbf{Sauvegarde incomplète en mode HA} : voir la
        section~\ref{sec:sauvegarde}.
  \item \textbf{Remise à zéro dangereuse} : \code{nettoyage_odoo.yaml} avec
        \code{full_reset=true} supprime \code{/var/lib/odoo}, devenu le point
        de montage du filestore partagé (section~\ref{sec:playbooks}).
\end{itemize}

% ===========================================================================
\section{Conclusion}
% ===========================================================================

Le projet fournit trois instances Odoo~14 qui partagent une base unique et un
filestore commun, avec une base PostgreSQL dont le primaire est remplacé
automatiquement en cas de panne. Les chemins critiques sont protégés : TLS de
bout en bout avec certificat client, autorités de certification séparées pour
SQL et etcd, montage du filestore exigé avant tout démarrage d'Odoo, et
initialisation unique de la base.

Trois chantiers restent prioritaires : rendre la sauvegarde opérationnelle en
mode HA (playbook et certificat client), mesurer les délais réels de bascule
et valider une restauration complète. Un répartiteur HTTP et la mobilité du
cron d'Odoo étendraient ensuite la disponibilité à la couche applicative.

% ===========================================================================
\appendix
\section{Catalogue des schémas}
% ===========================================================================

{\raggedright
Les schémas se trouvent dans le répertoire \code{docs/schemas-patroni-odoo14/}
du dépôt. Chacun existe en PNG (4800~$\times$~3200) pour la lecture, en SVG
pour l'impression et en draw.io pour la modification. Le fichier
\code{dossier-patroni-odoo14.drawio} les rassemble en quatorze onglets.\par}

\begin{table}[H]
\centering\small
\begin{tabularx}{\textwidth}{P{0.8cm} L P{6.4cm} P{0.9cm}}
\toprule
\rowcolor{gray!15}
\textbf{Nº} & \textbf{Schéma} & \textbf{Fichier (.png, .svg, .drawio)} & \textbf{Page} \\
\midrule
1  & Architecture globale : Odoo et PostgreSQL HA        & \code{01-architecture-globale-ha}        & \pageref{fig:01-architecture-globale-ha} \\
2  & Déploiement : hôtes Debian et conteneurs LXC        & \code{02-deploiement-hotes-lxc}          & \pageref{fig:02-deploiement-hotes-lxc} \\
3  & Patroni orchestre le rôle de chaque PostgreSQL      & \code{03-orchestration-patroni}          & \pageref{fig:03-orchestration-patroni} \\
4  & Connexion Odoo : du venv au primaire courant        & \code{04-connexion-odoo-haproxy}         & \pageref{fig:04-connexion-odoo-haproxy} \\
5  & Réplication : un primaire et deux répliques         & \code{05-replication-postgresql}         & \pageref{fig:05-replication-postgresql} \\
6  & Bascule automatique après panne du primaire         & \code{06-bascule-automatique}            & \pageref{fig:06-bascule-automatique} \\
7  & Quorum etcd et partitions réseau                    & \code{07-quorum-et-partitions}           & \pageref{fig:07-quorum-et-partitions} \\
8  & Bascule planifiée pour maintenance                  & \code{08-bascule-planifiee}              & \pageref{fig:08-bascule-planifiee} \\
9  & Réintégrer un ancien primaire sous Patroni          & \code{09-retour-ancien-primaire}         & \pageref{fig:09-retour-ancien-primaire} \\
10 & Transition : du cluster initial au mode Patroni     & \code{10-transition-vers-patroni}        & \pageref{fig:10-transition-vers-patroni} \\
11 & Flux réseau, certificats et surfaces de gestion     & \code{11-flux-reseau-et-tls}             & \pageref{fig:11-flux-reseau-et-tls} \\
12 & Sauvegarder puis restaurer Odoo et ses données      & \code{12-sauvegarde-et-restauration}     & \pageref{fig:12-sauvegarde-et-restauration} \\
13 & Installer Odoo 14 sur Debian dans les LXC           & \code{13-installation-odoo14-debian-lxc} & \pageref{fig:13-installation-odoo14-debian-lxc} \\
14 & Pile d'exécution : LXC, venv et service Odoo        & \code{14-pile-execution-odoo14-venv}     & \pageref{fig:14-pile-execution-odoo14-venv} \\
\bottomrule
\end{tabularx}
\caption{Catalogue des schémas de l'édition actuelle}
\end{table}

Les schémas des éditions précédentes se trouvent dans trois autres
répertoires de \code{docs/}, chacun avec son README, ses SVG et son fichier
draw.io.

\begin{table}[H]
\centering\small
\begin{tabularx}{\textwidth}{P{0.8cm} L P{7.4cm} P{0.9cm}}
\toprule
\rowcolor{gray!15}
\textbf{Nº} & \textbf{Schéma} & \textbf{Dossier de \texttt{docs/}, puis fichier (.png, .svg)} & \textbf{Page} \\
\midrule
15 & Raccordement IPVLAN              & \code{schemas-lxc/}\newline\code{01-raccordement-ipvlan}                          & \pageref{fig:schemas-lxc/01-raccordement-ipvlan} \\
16 & Accès utilisateur à Odoo         & \code{schemas-lxc/}\newline\code{02-acces-utilisateur-odoo}                       & \pageref{fig:schemas-lxc/02-acces-utilisateur-odoo} \\
17 & Flux Odoo vers PostgreSQL        & \code{schemas-lxc/}\newline\code{03-flux-odoo-postgresql}                         & \pageref{fig:schemas-lxc/03-flux-odoo-postgresql} \\
18 & Volume GlusterFS partagé         & \code{schemas-glusterfs/}\newline\code{01-volume-partage-trois-odoo}              & \pageref{fig:schemas-glusterfs/01-volume-partage-trois-odoo} \\
19 & Montage sur l'hôte et bind LXC   & \code{schemas-glusterfs/}\newline\code{02-montage-hote-et-bind-lxc}               & \pageref{fig:schemas-glusterfs/02-montage-hote-et-bind-lxc} \\
20 & GlusterFS et sauvegarde          & \code{schemas-glusterfs/}\newline\code{03-glusterfs-et-sauvegarde}                & \pageref{fig:schemas-glusterfs/03-glusterfs-et-sauvegarde} \\
21 & Vue de déploiement et flux       & \code{schemas-architecture-ha/}\newline\code{01-architecture-actuelle}            & \pageref{fig:schemas-architecture-ha/01-architecture-actuelle} \\
22 & Réplication et bascule manuelle  & \code{schemas-architecture-ha/}\newline\code{02-replication-et-bascule-manuelle}  & \pageref{fig:schemas-architecture-ha/02-replication-et-bascule-manuelle} \\
23 & Option HA Patroni, etcd, HAProxy & \code{schemas-architecture-ha/}\newline\code{03-option-ha-patroni-etcd-haproxy}   & \pageref{fig:schemas-architecture-ha/03-option-ha-patroni-etcd-haproxy} \\
24 & Bascule HA et quorum             & \code{schemas-architecture-ha/}\newline\code{04-bascule-ha-et-quorum}             & \pageref{fig:schemas-architecture-ha/04-bascule-ha-et-quorum} \\
\bottomrule
\end{tabularx}
\caption{Catalogue des schémas des éditions précédentes}
\end{table}

\end{document}
